% Converted from S8_Numbers_and_language/reviewer_response_package_20261120_claude_v8/RESPONSE_R2_23_EN.md (2026-11-17)
We thank the reviewer. ThyroidXAgent reports TI-RADS categories on the five levels of the ACR TI-RADS (TR1--TR5). The category is predicted from the nodule image by the TI-RADS branch of the multitask nodule model, which was trained on the ACR TI-RADS levels of the Stanford thyroid ultrasound cine-clip dataset and on DDTI. DDTI uses an earlier TI-RADS with subcategories 4a--4c, which we merged into level 4, so ThyroidXAgent does not output subcategories. The reference reports of SMU-HMC and ZJH-TS were written in routine practice at the two hospitals, which use the Chinese TI-RADS (C-TIRADS); most give a single category from 1 to 5, and 7 of the 400 SMU-HMC and 10 of the 150 ZJH-TS reports give subcategories 4a--4c. Our analyses do not require the two systems to agree: ThyClinScore checks only whether the impression states a TI-RADS category and does not compare categories, and in the reader study categories 4--5 were taken to indicate a malignant or suspicious nodule, which corresponds to the suspicious levels of both systems. The ``TI-RADS 4b'' in Supplementary Fig.~\SuppFigCaseReview{} was written by the Qwen3.5 Plus baseline, not by ThyroidXAgent.\par
